% !PRINTED 459
\weberpage{459}

Wir bilden eine lineare Function der $\alpha, \beta, \gamma \ldots$
\[
\xi = x\alpha + y\beta + z\gamma + \cdots,
\]
und beachten, dass jede der Gleichungen (1) nicht nur eine,
sondern mehrere Wurzeln hat, deren Zahl gleich dem Grade der
Gleichung ist. Ist also $\alpha', \beta', \gamma' \ldots$ irgend eine von
$\alpha, \beta, \gamma \ldots$ verschiedene Combination von je einer Wurzel von jeder der Gleichungen (1), so setzen wir:
\[
\xi' = x\alpha' + y\beta' + z\gamma' + \ldots
\]
und bilden auf die gleiche Weise $\xi'', \xi''' \ldots$. Die Anzahl der so
gebildeten Grössen $\xi$ ist, wenn $a$ der Grad von $A(x)$, $b$ der
Grad von $B(x)$, $c$ der Grad von $C(x)$ ist u.\,s.\,f.,
\[
m = abc \ldots
\]
Die Differenzen $\xi - \xi'$, $\xi - \xi''$, $\xi' - \xi'' \ldots$ sind lineare
Functionen von $x, y, z \ldots$ und keine von ihnen ist identisch
Null, da wir angenommen haben, dass keine der Gleichungen (1)
gleiche Wurzeln habe. Nach dem Satz 1 können wir also für
$x, y, z \ldots$ solche \gesperrt{rationale Zahlen} setzen, dass keine von
diesen Differenzen verschwindet, dass also die $m$ Werthe $\xi, \xi', \xi'' \ldots$
alle von einander verschieden sind.

Nun ist jede rationale Function, die in Bezug auf die Wurzeln
jeder der Gleichungen (1) symmetrisch ist, rational durch die
Co\"efficienten dieser Gleichungen ausdrückbar, d.\,h. es ist eine
Zahl in $\Omega$. Dazu gehören auch die Co\"efficienten der Function
$m^{\text{ten}}$ Grades von $t$:
\weq{2}{F(t) = (t - \xi)(t - \xi')(t - \xi'') \ldots,}
und $F(t) = 0$ ist also eine Gleichung in $\Omega$, die keine gleichen
Wurzeln hat, und deren eine Wurzel $\xi$ ist.

Ist ferner $\Theta$ eine Grösse in $\Omega(\alpha, \beta, \gamma \ldots)$, also eine ganze
rationale Function von $\alpha, \beta, \gamma \ldots$, und bezeichnen wir mit
$\Theta', \Theta'' \ldots$ die Grössen, die aus dieser rationalen Function
hervorgehen, wenn die Variablen durch $\alpha', \beta', \gamma' \ldots$ oder
$\alpha'', \beta'', \gamma'' \ldots$ u.\,s.\,f. ersetzt werden, so ist
\[
F(t) \left( \frac{\Theta}{t - \xi} + \frac{\Theta'}{t - \xi'} + \frac{\Theta''}{t - \xi''} + \cdots \right) = \psi(t)
\]
als symmetrische Function der $\alpha$ und der $\beta \ldots$ eine ganze
rationale Function $(m-1)^{\text{ten}}$ Grades von $t$ in $\Omega$, und wenn
man $t = \xi$ setzt, so folgt
%CONT>>
